\subsection{Diffusion Model for XRD Pattern Analysis}

This research employs a Denoising Diffusion Probabilistic Model (DDPM) \cite{ho2020denoising}
as a conditional denoiser that transforms ideal Rietveld-derived XRD patterns into realistic approximations of laboratory measurements.
The model learns the mapping from computed reference patterns to experimental data by conditioning on the dynamic time warping (DTW) distance between each synthetic--real pair, which serves as a proxy for the degree of structural and instrumental divergence.

\subsubsection{Model Architecture}

The implemented DDPM consists of two main phases: a forward diffusion process and a reverse denoising process similar to other diffusion models.
In the forward diffusion phase, synthetic XRD patterns undergo controlled noise addition, including Gaussian noise and application-specific noise.
This specific noise mimics the data augmentation patterns used in autoXRD \cite{szymanski2021probabilistic}, including peak shifting, peak intensity variations, peak removal, and increasing peak broadening. \\

The diffusion process implemented follows a denoising diffusion probabilistic model (DDPM) adapted for X-ray diffraction (XRD) data.
It introduces controlled noise to the data through a series of diffusion steps, progressively distorting the data from a clean state to a fully noisy state and subsequently learning to reverse this process.\\

\subsubsection{Augmentation Strategy}

The following section examines application-specific augmentation in detail, beginning with peak broadening of XRD peaks.
Peak broadening due to crystallite size reduction is modelled using Scherrer's equation \cite{scherrer1918determination}:
\begin{equation}
\text{FWHM} = \frac{K \cdot \lambda}{L \cdot \cos(\theta)}
\end{equation}
Where:
\begin{itemize}
    \item \( \text{FWHM} \) is the full width at half maximum (in radians).
    \item \( K \) is the shape factor (Scherrer constant, typically \(0.9\)).
    \item \( \lambda \) is the X-ray wavelength (default \(1.54056\) Å, Cu $K\alpha$ radiation).
    \item \( L \) is the crystallite size.
    \item \( \theta \) is the Bragg angle (half of the measured \(2\theta\)).
\end{itemize}

The crystallite size \( L \) varies linearly with the diffusion timestep fraction \( t_{\text{frac}} \) as:
\begin{equation}
L(t) = L_{\text{max}} - t_{\text{frac}}(L_{\text{max}} - L_{\text{min}})
\end{equation}
with \(L_{\text{min}}\) and \(L_{\text{max}}\) being minimum and maximum crystallite sizes, respectively.


Peak broadening is implemented via convolution with a Gaussian kernel whose width \( \sigma \) progressively increases as diffusion progresses:
\begin{equation}
\sigma = \sigma_{\text{base}}(t) \cdot \frac{L}{100}
\end{equation}
with \( \sigma_{\text{base}}(t) \) scaling from 0.02 (minimal broadening) to 1.0 (significant broadening) as:
\begin{equation}
\sigma_{\text{base}}(t) = 0.02 + t_{\text{frac}} \cdot 0.98
\end{equation}
The Gaussian kernel is:
\begin{equation}
k(x) = \frac{\exp\left(-\frac{1}{2}\left(\frac{x}{\sigma}\right)^2\right)}{\sum \exp\left(-\frac{1}{2}\left(\frac{x}{\sigma}\right)^2\right)}
\end{equation}


Additional augmentations introduced include:

\begin{itemize}
    \item \textbf{Peak Shifting}: Random shifting of spectra along the axis, increasing with timestep:
    \begin{equation}
    P_{\text{shift}}(t) = p_{\text{shift}} \cdot (1 + 2 \cdot t_{\text{frac}})
    \end{equation}

    \item \textbf{Peak Variations}: Random scaling of peak intensities within a dynamic range:
    \begin{equation}
    P_{\text{var}}(t) = p_{\text{var}} \cdot (1 + 2 \cdot t_{\text{frac}})
    \end{equation}
    Variation range:
    \begin{equation}
    \text{variation}_{\text{range}}(t) = \left[0.9 - 0.1 t_{\text{frac}}, 1.1 + 0.1 t_{\text{frac}}\right]
    \end{equation}

    \item \textbf{Removing Peaks}: Randomly zeroing peaks with probability:
    \begin{equation}
    P_{\text{remove}}(t) = p_{\text{remove}} \cdot (1 + 2 \cdot t_{\text{frac}})
    \end{equation}
    with the removal probability capped at 0.5.
\end{itemize}

\subsubsection{Forward Diffusion Process}
The forward diffusion introduces noise at timestep \( t \) as:
\begin{equation}
x_t = \sqrt{\bar{\alpha}_t} \cdot x_{0}^{\text{aug}} + \sqrt{1 - \bar{\alpha}_t} \cdot \epsilon
\end{equation}
where:
\begin{itemize}
    \item \( \bar{\alpha}_t = \prod_{s=1}^{t}(1 - \beta_s) \) is the cumulative product of alphas.
    \item \( x_{0}^{\text{aug}} \) is the augmented original data.
    \item \( \epsilon \sim \mathcal{N}(0, I) \) is Gaussian noise.
\end{itemize}

This formulation provides a physically motivated framework for synthesising and analysing X-ray diffraction spectra under realistic, noise-varied conditions, improving the model's ability to generalise across different measurement qualities.



\subsubsection{U-Net Noise Predictor}

The reverse process uses a 1D U-Net architecture that predicts the noise component added during the forward pass.
The network receives sinusoidal time embeddings concatenated with a DTW-based temperature conditioning signal, both projected to a shared hidden dimension of 16 channels through separate feed-forward networks with SiLU activations.

The encoder comprises two downsampling levels (channels 16$\to$32$\to$64) with stride-2 convolutions, each level containing two stochastic residual blocks with GroupNorm and SiLU non-linearities.
Self-attention layers are applied at levels one and two, capturing long-range correlations between distant peaks.
The decoder mirrors the encoder with transposed convolutions and skip connections, and a final sigmoid activation constrains the output to [0,\,1].
The full model has approximately 237\,000 trainable parameters.

\subsubsection{Training Procedure}


Two primary noise schedules are supported:
\begin{itemize}
    \item \textbf{Linear schedule}: Noise variance increases linearly from an initial value \( \beta_{\text{start}} \) to \( \beta_{\text{end}} \).
    \item \textbf{Cosine schedule} (default): Noise variance follows a cosine curve, providing smoother transitions between timesteps.
\end{itemize}

The cosine schedule is computed as:
\begin{equation}
\alpha_{\text{cumprod}}(t) = \frac{\cos^2\left(\frac{\frac{t}{T} + s}{1 + s} \cdot \frac{\pi}{2}\right)}{\cos^2\left(\frac{s}{1 + s} \cdot \frac{\pi}{2}\right)}
\end{equation}
\begin{equation}
\beta_t = 1 - \frac{\alpha_{\text{cumprod}}(t)}{\alpha_{\text{cumprod}}(t-1)}
\end{equation}
where \( T \) is the total number of timesteps and \( s = 0.008 \) is a small offset parameter.



The model training involves multiple steps:

\begin{itemize}
    \item \textbf{Synthetic Data Diffusion:} Synthetic patterns derived from CIF files are subjected to a stochastic diffusion process.
    Noise is gradually introduced to generate multiple noisy representations, improving the model's ability to generalise.
    \item \textbf{Denoising Training:} A U-Net model is trained to reverse the diffusion process,
    minimising the reconstruction error between original and denoised synthetic patterns.
    \item \textbf{Fine-Tuning:} To bridge the gap between theoretical and experimental data,
    the model undergoes fine-tuning using experimentally obtained XRD patterns,
    allowing adaptation to real-world variations and noise characteristics.
\end{itemize}

\subsubsection{Inference}

At inference the trained U-Net predicts the noise component $\hat{\epsilon}$ from a noisy input $x_t$.
The clean signal is then recovered analytically:
\begin{equation}
\hat{x}_0 = \frac{x_t - \sqrt{1-\bar{\alpha}_t}\,\hat{\epsilon}}{\sqrt{\bar{\alpha}_t}}
\label{eq:denoise}
\end{equation}
When $t=0$ the model acts as a direct mapping from ideal to realistic patterns, since the forward process contributes negligible noise at the first timestep.
To fit within a single 6\,GB GPU, inference is performed in mini-batches of 64 samples.

\subsubsection{Stochasticity and Variation Analysis}

A key feature of the DDPM is its stochastic nature: for the same ideal input, different random noise draws in the forward process produce different denoised outputs, each representing a plausible realisation of the measured pattern.
To quantify this variation we fix a single ideal input and apply the forward process at a small timestep ($t=50$) with $N=1\,000$ independent noise draws, denoising each one via Equation~\ref{eq:denoise}.
The resulting ensemble of outputs is analysed for its mean, standard deviation band, and peak-position coefficient of variation (CV), with a threshold of CV~$<$~5\,\% defining acceptable positional stability.
